%% kapa: exam=1 round=20 year=2009 period=1 form=A subject=civil_law subject_ko=민법 seq=30 number=30
%% source: Q-Net 감정평가사 기출문제 (cst00309) artlSeq=5024178 file=extracted/1008/5024178/2009년 제20회 감정평가사 1차시험 기출문제/1교시.hwp
%% conversion: kordoc 4.24.51 -> tools/gen1.py (auto); review=proofread-sonnet; flags=none
\begin{question}{30}
지역권에 관한 설명으로 옳지 \uline{않은} 것은?
\begin{choices}
  \item 지역권은 당사자 간의 계약에 의하여 성립할 수 있으나, 상린관계는 민법의 규정에 의하여 당연히 인정된다.
  \item 지역권은 요역지를 위한 종된 권리이므로 요역지와 분리하여 양도하지 못한다.
  \item 지역권은 계속되고 표현된 것에 한하여 시효취득의 대상이 될 수 있다.
  \item 지역권이 설정된 후에 요역지에 대하여 지상권을 취득한 자는 특약이 없는 한 그 지역권을 행사할 수 없다.
  \item 요역지와 승역지는 반드시 인접하고 있을 필요는 없다.
\end{choices}
\end{question}
\answer{4}
